\documentclass{book}

\usepackage{graphicx,epsfig}
\usepackage{hyperref}
\usepackage{amsmath,amssymb}
\usepackage[toc,page]{appendix}


\renewcommand{\baselinestretch}{1}

\setlength{\textheight}{9in}
\setlength{\textwidth}{6.5in}
\setlength{\headheight}{0in}
\setlength{\headsep}{0in}
\setlength{\topmargin}{0in}
\setlength{\oddsidemargin}{0in}
\setlength{\evensidemargin}{0in}
\setlength{\parindent}{.3in}
\pagestyle{plain}





\begin{document}
\bibliographystyle{plain}

\large

\thispagestyle{empty}
\centerline{\bf A REPORT}
\vspace*{0.3cm}
\centerline{\bf ON}
\vspace*{0.3cm}
\centerline{\bf AI-POWERED EXECUTIVE DASHBOARD \& BUSINESS REVIEW AUTOMATION} % enter here
\vspace*{2cm}

\centerline{BY}
\vspace*{1cm}

\begin{center}
	\begin{tabular}{lll} % enter here
		Name of the student & \hspace*{1cm} ID No. & \hspace*{1cm} Discipline \\ 
		Sambhav Jain &\hspace*{1cm} 2024B4A80840H &\hspace*{1cm} M.Sc. Math \& B.E. ENI \\
	\end{tabular}
\end{center}

\vspace*{2cm}
\begin{center}
	{\bf Prepared in partial fulfillment of the \\
	Practice School - I} \\ % enter here
\end{center}

\vspace*{0.5cm}

\centerline{AT}




% \begin{document}
% \bibliographystyle{plain}

% \large

% \thispagestyle{empty}
% \centerline{\bf A REPORT}
% \vspace*{0.3cm}
% \centerline{\bf ON}
% \vspace*{0.3cm}
% \centerline{\bf SUPPORT VECTOR MACHINES} % enter here
% \vspace*{2cm}

% \centerline{BY}
% \vspace*{1cm}

% \begin{center}
% 	\begin{tabular}{lll} % enter here
% 		Name(s) of the students & \hspace*{2cm} ID No(s). & \hspace*{2cm} Discipline(s) \\ 
% 		Amardeep Singh Grewal &\hspace*{2cm} 2022A2PS0851P &\hspace*{2cm} B. E in Civil Engineering \\
% 		Harish Kumar &\hspace*{2cm} 2021B4PS0642H &\hspace*{2cm} M. Sc. in Mathematics \\
% 		Akhsay Agrawal & \hspace*{2cm} 2022A2PS1030P & \hspace*{2cm} B. E. in Civil Engineering\\
% 	\end{tabular}
% \end{center}

% \vspace*{2cm}
% \begin{center}
% 	{\bf Prepared in partial fulfillment of the \\
% 	Practice School - II} \\ % enter here
% \end{center}

% \vspace*{0.5cm}

% \centerline{AT}






\vspace*{1cm}
        \centerline{\bf Celebal Technologies, Jaipur} % enter here
	\vspace{0.2cm}
	\centerline{\bf A Practice School - I}
	\vspace{0.2cm}
	\centerline{\bf Station of}
	\vspace{0.5cm}
\begin{figure}[ht]
\epsfxsize=1.0in
\centerline{\epsffile{logo.eps}}
\end{figure}

	\centerline{\bf BIRLA INSTITUTE OF TECHNOLOGY \& SCIENCE, PILANI}
	\vspace*{0.5cm}
	\centerline{\bf June, 2026} % enter here
	\newpage

\thispagestyle{empty}

	\centerline{\bf BIRLA INSTITUTE OF TECHNOLOGY \& SCIENCE, PILANI}
	\vspace*{0.2cm}
	
    
    \centerline{\bf Practice School Division}
\vspace*{0.5cm}

\begin{center}
    \begin{tabular}{@{}p{5.5cm} p{9.5cm}@{}}
    
    Station: & Celebal Technologies \\[0.5cm]
    
    Centre: & Jaipur \\[0.5cm]
    
    Duration: & 8 Weeks \\[0.5cm]
    
    Date of start: & 25 May 2026 \\[0.5cm]
    
    Date of submission: & 18 July 2026 \\[0.5cm]
    
    Title of the project: & AI-Powered Executive Dashboard \& Business Review Automation for the CEO Office \\[0.5cm]
    
    ID No.(s), Name(s) and Disciplines of the student: & 2024B4A80840H, Sambhav Jain, M.Sc. Mathematics \& B.E. Electronics and Instrumentation \\[0.5cm]
    
    Name(s) and Designation(s) of expert(s): & [Mentor Name], [Mentor Designation] \\[0.5cm]
    
    Name(s) of PS faculty member(s): & Prof. S Gangopadhyay \\[0.5cm]
    
    Key words: & Python, Pandas, Streamlit, LLMs, Dashboarding, ETL \\[0.5cm]
    
    Project area(s): & Data Engineering, Analytics, Full-Stack Web Development \\[0.5cm]
    
    Abstract: & This project focuses on building an automated reporting system that provides real-time visibility into company-wide KPIs including revenue, project health, headcount, and client satisfaction. Data is collected from multiple sources, cleaned and integrated using Python and Pandas, and visualized via an interactive Streamlit dashboard. Additionally, an LLM-powered briefing generator is integrated to automatically summarize key business highlights, risks, and strategic action items for executive leadership. \\
    			& \\
			& \\
			& \\
			& \\
			Signature of the student(s) & \hspace*{1cm}Signature of the PS faculty member(s) \\
			& \\
			Date & \hspace*{1cm}Date \\ % enter here
    \end{tabular}
\end{center}




    % \centerline{\bf Practice School Division}
	% \vspace*{0.3cm}
	% {\bf 
	% %\begin{center}
	% 	\begin{tabular}{p{6cm}l}
	% 		& \\
	% 		& \\
	% 		& \\
	% 		Station & \hspace*{3cm} Centre \\ Celebal Technologies & \hspace*{3cm} Jaipur \\
	% 		& \\
	% 		Duration &\hspace*{3cm} Date of start \\ % enter here
	% 		& \\
	% 		Date of submission & {} \\ % enter here
	% 		& \\
	% 		Title of the project: & \\ % enter here
	% 		& \\
	% 		ID No.(s), Name(s) and Disciplines of the student(s): & \\
	% 		& \\ % enter here
	% 		& \\ % enter here
	% 		& \\ % enter here
	% 		Name(s) and Designation(s) of the expert(s): & \\
	% 		& \\ % enter here
	% 		& \\ % enter here
	% 		Name(s) of the PS faculty member(s): & \\
	% 		& \\ % enter here
	% 		Key words: & \\ % enter here
	% 		& \\
	% 		Project area(s):  & \\ % enter here
	% 		& \\
	% 		Abstract: & \\ % enter here
	% 		& \\
	% 		& \\
	% 		& \\
	% 		& \\
	% 		& \\
	% 		Signature of the student(s) & \hspace*{1cm}Signature of the PS faculty member(s) \\
	% 		& \\
	% 		Date & \hspace*{1cm}Date \\ % enter here
	% 	\end{tabular}}



\newpage
\centerline{\Large \bf Acknowledgements}
\vspace*{1cm}

\par\noindent
I would like to express my sincere gratitude to the Practice School Division of BITS Pilani and my host organization, Celebal Technologies, Jaipur, for providing me with the incredible opportunity to work on the ``AI-Powered Executive Dashboard \& Business Review Automation'' project. This Practice School-I experience has been highly instrumental in bridging the gap between academic theory and real-world industry application.

\vspace*{0.4cm}
\par\noindent
I am deeply thankful to my Practice School faculty in-charge, Prof. S Gangopadhyay, for his continuous guidance, encouragement, and valuable evaluation throughout the duration of this project. His academic insights helped shape the structural and professional direction of my work.

\vspace*{0.4cm}
\par\noindent
I would also like to extend my heartfelt appreciation to my industry mentor, {[Mentor Name]}, {[Mentor Designation]} at Celebal Technologies. Their practical expertise, patience, and technical mentorship in Python, data engineering, and Streamlit were crucial in helping me successfully build and deploy the final dashboard application.

\vspace*{0.4cm}
\par\noindent
Finally, I would like to thank my peers for their support during this entire eight-week journey.

	\vfill

		\par\noindent {\bf Signature of the student(s)}




	% \newpage
	% \centerline{\bf \Large Acknowledgements}
	% \vspace*{0.5cm}
	% 	\par\noindent Ensure that you start each paragraph with \begin{verbatim} \par\noindent \end{verbatim} so that the initial spacing is not there in the text.  % enter here
		
	% 	\href{https://www.example.com}{Visit our website} 
	% 	\vfill

	% 	\par\noindent {\bf Signature of the student(s)}



\tableofcontents
\listoffigures
\listoftables







\chapter{Introduction}

\section{Background and Motivation}
In the modern corporate landscape, executive leadership teams, particularly the Chief Executive Officer (CEO), require immediate and accurate visibility into organizational performance. Decision-making at the highest level relies heavily on understanding the real-time health of various business units. However, as organizations grow, the sheer volume of operational data makes manual reporting inefficient and prone to human error. There is a critical need for automated systems that can instantly consolidate data and present actionable insights.

\section{Problem Statement}
Currently, enterprise data is highly fragmented across disparate systems. Customer Relationship Management (CRM) tools track sales pipelines and client satisfaction, Human Resource Management Systems (HRMS) monitor employee headcount, and Finance ledgers record monthly revenues and operational costs. For a CEO, manually extracting, cleaning, and correlating data across these distinct platforms to prepare for a weekly business review is a time-consuming bottleneck that delays strategic action.

\section{Proposed Solution}
This project addresses this bottleneck by developing an ``AI-Powered Executive Dashboard \& Business Review Automation'' system. By leveraging Python-based data engineering libraries, the system automatically extracts and processes raw data from multiple operational spreadsheets. This cleaned data is then fed into an interactive, real-time frontend dashboard. Furthermore, the system integrates Generative AI to autonomously analyze the metrics and produce a written, executive-grade strategic briefing, highlighting critical risks and financial highlights.

\section{Objectives}
The primary objectives of this Practice School project at Celebal Technologies are as follows:
\begin{itemize}
    \item \textbf{Data Pipeline Engineering:} To build a robust Extract, Transform, Load (ETL) pipeline using the Pandas library capable of merging disparate CRM, HRMS, and Finance datasets while programmatically handling missing or dirty data.
    \item \textbf{Interactive Visualization:} To develop a responsive, dark-theme frontend analytics dashboard using the Streamlit framework to visualize Key Performance Indicators (KPIs) through interactive charts and metric cards.
    \item \textbf{AI Integration:} To integrate Google's Gemini Large Language Model (LLM) via API to autonomously analyze live operational data and generate structured strategic action items.
    \item \textbf{Workflow Automation:} To combine these components into a single, executable application that simulates a seamless weekly data refresh for the CEO office.
\end{itemize}










% \chapter{Literature review}
% You may put in all the relevant papers and books that you have referred and also give references to them by citing them appropriately. Ensure that the bibtex entries are updated to reflect them correctly. 
% \chapter{Methodology}
% Specify the requirements in the project and the methods you wish to use. A brief justification on the method would be useful. 



\chapter{Literature Review}

\section{The Evolution of Business Intelligence}
Historically, enterprise reporting and Business Intelligence (BI) relied heavily on static spreadsheet software or incredibly complex, monolithic BI tools that required specialized IT teams to operate. While these legacy systems provided value, they suffered from a critical flaw: latency. By the time data was extracted, cleaned, and visualized for a Chief Executive Officer (CEO), the insights were often days or weeks out of date. Modern literature on organizational agility emphasizes the need for real-time, programmatic data pipelines that can seamlessly transform raw data into actionable insights without human intervention \cite{mckinsey_agile_2020}.

\section{Data Engineering with Python and Pandas}
Python has emerged as the lingua franca of data engineering, largely due to the power of the Pandas library. Pandas introduces the `DataFrame`, a highly efficient, in-memory data structure that allows developers to manipulate vast datasets using vectorized operations. Unlike traditional row-by-row iteration found in older programming paradigms, Pandas leverages underlying C-language optimizations to perform aggregations, joins, and missing-value imputations in fractions of a second \cite{mckinney2010data}. For an executive dashboard, where multiple disparate data sources (HRMS, CRM, Finance) must be merged continuously, Pandas serves as the perfect extraction and transformation engine.

\section{Streamlit: Democratizing Web Applications}
Traditionally, building a web-based dashboard required a full-stack engineering team proficient in HTML, CSS, JavaScript, and backend frameworks like Django or Flask. Streamlit fundamentally disrupted this paradigm. It allows data scientists and engineers to build highly interactive, responsive web applications using only pure Python \cite{streamlit_docs}. By abstracting away the complexities of frontend routing and state management, Streamlit enables rapid prototyping. This literature heavily influenced the decision to use Streamlit for this project, as it allowed maximum focus on data logic and UI/UX design rather than boilerplate web development.

\section{Generative AI in Enterprise Contexts}
The most recent paradigm shift in business analytics is the integration of Large Language Models (LLMs). While traditional dashboards can show a CEO *what* the data is (e.g., a drop in revenue), they struggle to explain *why* it matters or *what* to do next. Modern APIs, such as Google's Gemini, possess the semantic reasoning capabilities to ingest structured numeric data and output nuanced, human-readable strategic summaries \cite{gemini_api}. Integrating LLMs directly into BI pipelines represents the cutting edge of automated executive reporting.


\chapter{Methodology}

The development of the AI-Powered Executive Dashboard was executed in four distinct, iterative phases. This approach ensured that each layer of the application—from the raw data to the final user interface—was thoroughly tested and optimized.

\section{Phase 1: Synthetic Data Engineering}
At the inception of the project, direct access to live, sensitive corporate enterprise databases was restricted due to standard security protocols. Rather than halting development, a Python-based synthetic data generator was engineered. 
Using the `numpy` and `pandas` libraries, an automated script was written to generate thousands of rows of realistic corporate data across three simulated databases:
\begin{itemize}
    \item \textbf{CRM Pipeline:} Simulating client acquisitions, project health (categorized as Red, Yellow, or Green), and satisfaction scores.
    \item \textbf{HRMS Ledger:} Tracking employee headcount, attrition rates, and departmental allocations.
    \item \textbf{Financial Records:} Generating monthly revenue streams, operational costs, and profit margins.
\end{itemize}
To simulate the imperfections of real-world databases, the script intentionally injected `NaN` (Not a Number) values and formatting inconsistencies, ensuring the downstream pipeline would be robust enough to handle "dirty" data.

\section{Phase 2: The ETL Pipeline Construction}


With the raw data generated, the next step was building the Extract, Transform, Load (ETL) architecture. A dedicated 'data\_processor.py' module was developed to serve as the backend engine. 

\par\noindent
First, the script extracted the raw CSV files. During the Transformation phase, advanced Pandas functions were utilized to clean the data. Missing numeric values were programmatically replaced using median imputation to prevent skewed averages, while missing categorical data was flagged for review. The disparate datasets were then joined together using a common 'Project\_ID' primary key. Finally, the script calculated derived business metrics, such as the Net Profit Margin, which is a critical KPI for executive review.



\section{Phase 3: Streamlit Frontend Development}
The transformed data required an intuitive, frictionless user interface. Streamlit was utilized to build a dark-themed, single-page web application. The methodology for the UI/UX design was strictly tailored for executive consumption—prioritizing immediate clarity over visual clutter.
\par\noindent
At the top of the dashboard, large "Metric Cards" were implemented to display top-line numbers (Total Revenue, Active Headcount) alongside delta indicators (showing percentage growth from the previous quarter). Below the metric cards, interactive visualizations were rendered, allowing the user to filter the data by specific quarters or departments dynamically.

\section{Phase 4: LLM Integration and Prompt Engineering}
The final, most advanced phase of the methodology was the integration of the Generative AI briefing engine. The cleaned, aggregated data metrics were passed from the Pandas backend into a structured text prompt.
\par\noindent
Prompt engineering was a critical challenge here. The LLM was explicitly instructed to adopt the persona of a Senior Business Analyst. Through the API connection, the LLM ingests the weekly data, analyzes the correlations (e.g., noting if rising costs are misaligned with stagnant revenue), and generates a concise, three-bullet-point strategic briefing. This ensures the CEO is not just looking at charts, but is actively receiving automated, data-backed business advice.





% \chapter{Results and discussions}
% Collate all the results in this chapter and discuss why they are so. Ensure that the plots are visible to reader. Yuu may use pstricks package to enhance some aspects in a figure. 
% \chapter{Conclusions and future work}
% Interpret the results and draw conclusions based on the work done. You should also include a paragraph or two on what could be done to improve the results. 

% \chapter{Recommendations}
% Use this chapter to include all recommendations.


\chapter{Results and Discussions}

\section{ETL Pipeline Performance}
The first major result of this project was the successful execution of the automated ETL pipeline. The Python script (\texttt{data\_processor.py}) successfully ingested the disparate CRM, HRMS, and Finance CSV files. The Pandas library handled the intentionally injected missing \texttt{NaN} values via median imputation seamlessly. The processing time for merging and cleaning the synthetic datasets averaged under two seconds locally, demonstrating the high efficiency of Pandas' vectorized operations.

\section{Dashboard Interface and Usability}
The Streamlit frontend successfully translated the processed backend data into an interactive visual format. The dark-themed User Interface (UI) rendered without latency. Key Performance Indicators (KPIs) such as Total Revenue, Profit Margins, and Employee Attrition were prominently displayed on top-level metric cards.

% -------- SCREENSHOT PLACEHOLDER --------
\begin{figure}[h]
    \centering
    % STEP 1: Upload your screenshot to Overleaf.
    % STEP 2: Delete the \vspace line below.
    % STEP 3: Remove the % symbol from the \includegraphics line and put your image name inside the brackets.
    
    \vspace{6cm} 
    % \includegraphics[width=0.85\textwidth]{your_image_name.png}
    
    \caption{The interactive Streamlit Executive Dashboard displaying real-time KPIs.}
    \label{fig:dashboard1}
\end{figure}
% ----------------------------------------

As seen in Figure \ref{fig:dashboard1}, the interactive charts allow the end-user to hover over specific data points to view exact numerical values, a significant upgrade over static spreadsheet graphs.

\section{AI-Generated Strategic Briefings}
The most critical result was the successful integration of the Google Gemini API. Upon refreshing the dashboard, the LLM successfully ingested the aggregated JSON payload of the weekly metrics. Within seconds, the API returned a highly structured, executive-grade strategic briefing. The AI successfully identified complex correlations, such as flagging potential operational risks when HR headcount increased while CRM project acquisitions decreased, proving its value as an automated analytical tool.


\chapter{Conclusions and Future Work}

\section{Conclusions}
The ``AI-Powered Executive Dashboard \& Business Review Automation'' project successfully demonstrated that manual, time-intensive enterprise reporting can be fully automated. By combining the data-wrangling capabilities of Python, the rapid UI deployment of Streamlit, and the semantic reasoning of Large Language Models, the project delivered a prototype that provides executive leadership with real-time, actionable insights. This system fundamentally shifts the paradigm from executives merely looking at historical data to interacting with a real-time business advisor.

\section{Future Work}
While the current iteration operates successfully as a proof-of-concept, several enhancements are required before enterprise-grade production deployment:
\begin{itemize}
    \item \textbf{Live Database Integration:} Transitioning the data ingestion layer from local static CSV files to live cloud databases (such as PostgreSQL or MongoDB) via secure API endpoints.
    \item \textbf{Authentication and Role-Based Access:} Implementing robust OAuth 2.0 security protocols to ensure that highly sensitive CEO-level financial data is strictly access-controlled.
    \item \textbf{Cloud Deployment:} Migrating the application from a local host environment to a scalable cloud architecture using Docker containers on platforms like AWS or Microsoft Azure.
\end{itemize}


\chapter{Recommendations}
Based on the successful development of this prototype during the Practice School duration, the following recommendations are proposed:
\begin{itemize}
    \item \textbf{Client Demonstration:} It is recommended that Celebal Technologies utilize this dashboard architecture as a Minimum Viable Product (MVP) to demonstrate the power of LLM integration to prospective clients seeking Business Intelligence modernization.
    \item \textbf{Open-Source Alternatives:} While the Gemini API provided excellent results, evaluating open-source LLMs (such as Llama 3) for the briefing engine could reduce long-term API token costs and enhance data privacy, ensuring that sensitive corporate data does not need to be transmitted to external servers.
\end{itemize}




% \bibliography{ref}

% \addcontentsline{toc}{chapter}{Bibliography}

% \begin{appendices}
% \chapter{The proof of Central Limit Theorem}
% Show the proof here
% \chapter{Riemann Lebesgue Lemma}
% 	\label{appB}
% Prove the lemma here
% \end{appendices}


% \end{document}


\bibliography{ref}
\addcontentsline{toc}{chapter}{Bibliography}

\begin{appendices}
\chapter{Core Pipeline Code Snippets}
The following snippet demonstrates the Pandas logic used to clean the generated dataset and handle missing values within the ETL pipeline.

\begin{verbatim}
import pandas as pd
import numpy as np

def clean_enterprise_data(file_path):
    # Load raw mock data
    df = pd.read_csv(file_path)
    
    # Handle intentionally injected NaN values
    # Replace missing numerical values with the median
    df['Revenue'].fillna(df['Revenue'].median(), inplace=True)
    
    # Flag missing categorical data
    df['Project_Health'].fillna('NEEDS REVIEW', inplace=True)
    
    return df
\end{verbatim}

\end{appendices}

\end{document}